% !TeX root = manuscript.tex
% Author-provided replacement abstract; wording retained as requested.
\begin{abstract}
Automatic video dubbing requires intelligible speech that preserves a reference voice while remaining synchronized with visible articulation. Although implicit alignment avoids the external forced aligners required by duration-based approaches, early audiovisual fusion and unidirectional text conditioning can hinder reciprocal multimodal interaction and fine-grained temporal correspondence. To address these issues, we propose JointDub, a locality-aware multimodal diffusion transformer for synchronized video dubbing. Specifically, modality-specific streams separately process speech latents, visual features, and textual representations, and exchange information through joint attention before audio-only refinement. This design enables coupled representation learning while retaining modality-specific processing. Furthermore, we introduce Gaussian-windowed audiovisual attention, which applies a soft temporal prior to bidirectional audio-video interactions while preserving global access to textual context. By favoring temporally neighboring cross-modal associations, this mechanism promotes lip synchronization without imposing externally specified phoneme boundaries. The model is trained using conditional flow matching and connectionist temporal classification, without externally supplied phoneme durations. Experiments on two public datasets demonstrate state-of-the-art overall performance among publicly available dubbing models, considering speech intelligibility, speaker similarity, and audiovisual synchronization. Extensions to video-to-speech synthesis and visual forced alignment further demonstrate the versatility of the framework.
\end{abstract}
